\section{Migration plan}\label{sec:migration}

The refactor described in the preceding sections rewrites the very heart of the
composition math---the meaning of the stored child offset, the input to
\src{model/parts/Part.cpp}{169}'s two-pass mass/CM/parallel-axis loop, and the axial
station threaded through \src{model/parts/Part.cpp}{214}'s aero walk. A change of that
reach to a shipping simulator cannot be validated by inspection. The migration is
therefore structured around a single load-bearing claim: the new geometry-driven
placement model is \emph{physics-invariant} on every existing design. The center of
mass, the composite inertia tensor, and every resolved axial station that the legacy
CM-to-CM reader produces must be reproduced \emph{bit-for-bit} by the migrated reader.
The plan is consequently analytic rather than empirical---no flight re-runs, no
tolerance fudging---and it is gated by a test that is written and made to pass
\emph{before} any production type changes.

The four phases below are ordered so that the build stays green at every step and the
invariance claim is proven, not assumed.

\subsection{Phase 0 --- the invariance gate, written first}

Before a single type changes, a test loads every fixture in
\srcf{tests/data/designs} (the 24 \code{.qrd} designs spanning the \code{micro13}
through \code{xl75} ladders) through the \emph{current} reader and snapshots three
quantities per design: \cppinline|getCompositeCm(0)| (\src{model/parts/Part.h}{134}),
\cppinline|getCompositeI(0)| (\src{model/parts/Part.h}{146}), and the resolved axial
station of every part in the existing \zfwd{} CM-to-CM convention.

The discipline that makes this a \emph{gate} rather than a self-fulfilling check is that
the snapshot is taken from the legacy reader's own output, not re-derived from the
design's intended geometry. A re-derivation could silently share whatever bug the
legacy path contains; the legacy reader's literal output cannot. These frozen values
become the ground truth the migration is held against.

\begin{designrule}[Invariance is defined by the incumbent, not the spec]
The Phase-0 snapshot is whatever the current CM-to-CM code computes today, including
quantities the new model would consider mis-placed. The migration's job is to reproduce
the incumbent numerics exactly; \emph{correcting} a design's geometry (the
\code{xl75\_multi} poke-through) is deferred to Phase 3, where it is done deliberately
and visibly, not folded silently into the type swap.
\end{designrule}

\subsection{Phase 1 --- types and resolver land, both APIs coexist}

Phase 1 introduces the new machinery alongside the old without removing anything the
build depends on. The new header \srcf{model/parts/Placement.h}
(\cref{sec:datamodel}) lands with \seat{Abut}/\seat{NestInBore}/\seat{OnSurface}, the
\code{StationLink}, \code{Station}, and \code{Pose} value types, and the free-function
resolver and overlap sweep (\cref{sec:resolver}, \cref{sec:diagnostics}). The geometry
virtuals---\code{stationAt}, \code{radiusOuterAt}/\code{radiusInnerAt},
\code{axialLength}, \code{isSolid}, \code{cmStationLocal}---are added to the base
\code{Part} and overridden on the four concrete part types. The single child-storage
slot at \src{model/parts/Part.h}{326} is swapped in place from
\cppinline|std::tuple<std::shared_ptr<Part>, Vector3>| to
\cppinline|std::pair<std::shared_ptr<Part>, StationLink>|, and
\src{model/parts/Part.cpp}{169} and \src{model/parts/Part.cpp}{214} are rewritten to
consume the resolver's geometric pose.

Crucially, the legacy authoring surface survives. The current sole overload
\src{model/parts/Part.h}{239}, \cppinline|addChildPart(std::shared_ptr<Part>, Vector3)|,
is retained as a \cppinline|[[deprecated]]| shim so that every existing tuple-style call
site and every legacy \code{.qrd} fixture still compiles and still produces the
legacy placement. The shim is the analytic core of the migration: it must synthesize a
\code{StationLink} whose resolved geometry reproduces the legacy CM-to-CM offset
exactly.

\subsubsection{The CM-to-CM recovery formula}

The legacy \code{Vector3} is the child's CM expressed relative to the parent's CM. The
resolver, by contrast, places a child by its \emph{aft-plane origin} in the parent
frame (\cref{sec:resolver}). The shim bridges the two by inverting the CM-to-station
relationship, working entirely in the \zfwd{} frame so that no axis is mirrored.
\Cref{lst:shim} gives the recovery.

\begin{listing}[htbp]
\begin{minted}{cpp}
// All quantities in +z = forward. legacyOffset.z is the child CM minus the
// parent CM (the legacy CM-to-CM coupling).
//
// 1. Recover the child's aft-plane origin in the parent frame. *CmStationLocal
//    is cmStationLocal() evaluated in the +z-forward local frame; the cone is
//    the only part whose adapter differs (see the cone frame adapter).
childOriginZ_fwd = parentCmStationLocal - childCmStationLocal + legacyOffset.z;

// 2. Choose a station pair, then back the gap out of the resolved geometry.
//    childStation / parentStation are the picked Station landmarks.
gap_unsigned = childOriginZ_fwd + childStation.z - parentStation.z;  // unsigned

// 3. Infer the seat from the PART PAIR -- BEFORE applying any sign:
//      tube fore-rim against tube aft-rim         => Abut
//      small-tube OD inside large-tube ID         => NestInBore
//      fin root on the wall                       => OnSurface
seat = inferSeat(parent, child);

// 4. Apply the seat sign LAST. NestInBore inserts aft (-z); all others stand
//    off forward (+z).
gap = (seat == SeatKind::NestInBore) ? -gap_unsigned : +gap_unsigned;
\end{minted}
\caption{The \code{[[deprecated]]} shim's analytic recovery of a
\code{StationLink} from a legacy CM-to-CM \code{Vector3}. The \code{cmStationLocal}
adapter is consumed here (its first production use); the cone supplies the only
non-trivial override.}\label{lst:shim}
\end{listing}

Two details in \cref{lst:shim} carry the correctness of the entire migration, and both
are easy to get backwards.

\begin{keyinsight}[+z is forward throughout --- no mirror]
The legacy \code{xl75\_multi} coupler stores \cppinline|z = +0.49| relative to the body
(\src{tests/data/designs/xl75\_multi.qrd}{19}). Because the convention is \zfwd{} on both
sides of the recovery, that \cppinline|+0.49| is read directly as \emph{forward}---toward
the nose tip---with no sign flip and no whole-rocket inversion. The migration is a pure
re-expression of the same axis, which is precisely why \zfwd{} was chosen over \zaft{}
(\cref{sec:philosophy}): \zaft{} would have forced a sign-negating rewrite of every
fixture here. The cone is the sole part whose \emph{internal} CM frame
(\code{coneCmOffset}, \src{model/parts/ConicalNoseCone.cpp}{49}) opposes the forward
convention, and that single discrepancy is absorbed by the
\code{cmStationLocal} adapter (\cref{sec:datamodel}), never by mirroring the link.
\end{keyinsight}

\begin{keyinsight}[Seat is inferred before the gap sign]
The seat must be determined from the part pair \emph{before} the gap sign is applied,
because \seat{NestInBore} is the one seat that flips the sign (step~4 of
\cref{lst:shim}). Computing an unsigned gap first, then choosing the sign from a
seat inferred independently of that gap, is what makes the ``reproduces the exact
station'' invariant hold for \emph{every} seat---not only the sign-preserving
\seat{Abut} and \seat{OnSurface} cases. Inferring the seat from the already-signed gap
would be circular and would mis-place every nested coupler.
\end{keyinsight}

By the end of Phase 1 the build is green, the new and legacy APIs coexist, and the shim
reproduces legacy behavior analytically.

\subsection{Phase 2 --- run the gate against the migrated reader}

Phase 2 loads every fixture through the now \code{0.2}-capable reader---legacy
\cppinline|<offset>| elements flow through the shim of \cref{lst:shim} into
\code{StationLink}s---and asserts \emph{bit-identical} \cppinline|getCompositeCm(0)|,
\cppinline|getCompositeI(0)|, and resolved axial stations against the Phase-0
snapshots. Equality is exact, not within a tolerance: the migration is a re-expression
of the same arithmetic over the same geometry, so any deviation is a defect, and a
tolerance would merely hide it.

A dedicated sub-test pins the nose CM to a hand-computed value as the specific guard
that the cone frame adapter is correct. For the \code{xl75\_multi} nose---a solid cone
of length \cppinline|L = 0.30| (\src{tests/data/designs/xl75\_multi.qrd}{5})---the
centroid sits at $L/4$ forward of the base, i.e. \cppinline|cmStationLocal = 0.075|.
This is the forward-frame image of \code{coneCmOffset}'s internal $z = L/2 - \bar{h}$
with $\bar{h} = L/4$ (\src{model/parts/ConicalNoseCone.cpp}{49}); the adapter maps it to
$L/2 - (L/2 - \bar{h}) = \bar{h} = L/4$ (\cref{sec:datamodel}). Pinning this one scalar to a value computed
by hand---independently of both the legacy and migrated code paths---is the test that
catches a sign or frame error in the single most correctness-critical conversion in the
design.

\begin{keyinsight}[Why the cone gets its own hand-checked sub-test]
A symmetric part has its CM at mid-length, so \code{cmStationLocal} is $L/2$ and any
frame error cancels by symmetry. The cone is the only current part whose CM is
off-center, so it is the only part for which a forward-vs-internal frame mistake would
survive the bit-identical aggregate check yet still be wrong. The hand-computed
\cppinline|0.075| isolates exactly that risk.
\end{keyinsight}

\subsection{Phase 3 --- cut over the fixtures and retire the shim}

With invariance proven, Phase 3 re-saves the corpus to version \code{0.2}
(\cref{sec:serialization}), replacing each \cppinline|<offset>| with an explicit
\cppinline|<link>|. Here the migration deliberately stops preserving the past. The
\code{xl75\_multi} coupler---legacy-placed forward into and through the nose by its
\cppinline|+0.49| CM-to-CM offset---migrates to a \seat{NestInBore} link, and at that
point its Layer-2 envelope sweep \emph{fires}: the coupler outer radius exceeds the
tapering nose skin forward of the body fore rim, and the resolver returns
\cppinline|solve.ok == false| with a located \code{OverlapDiagnostic}
(\cref{sec:diagnostics}, \cref{sec:worked}).

\begin{hardfail}[The migration reports the pre-existing bug instead of preserving it]
The legacy file drew an interpenetrating coupler silently because the simulator and the
visualizer each consumed the bad offset in their own divergent way (\cref{sec:intro}).
Once the relationship is expressed as ``nest into the body's fore bore,''
\seat{NestInBore} constrains the insertion direction and the envelope sweep makes the
forward poke-through a hard, located \fail{}. The correct response is not to suppress
the diagnostic but to fix the fixture---reduce the insertion depth or shorten the
coupler---until it resolves \ok{}.
\end{hardfail}

After each fixture resolves clean, the invariance gate is re-run against the fresh
\code{0.2} files to confirm the corrected designs are stable under reload. Only then is
the \cppinline|[[deprecated]]| overload at \src{model/parts/Part.h}{239} removed,
collapsing the authoring surface onto the single \code{StationLink} API
(\cref{sec:authoring}) and completing the cutover. The accessor side---\code{getChildParts()}
returning a \code{pair} rather than a \code{tuple}---is migrated explicitly at its
enumerated call sites (\cref{sec:files}); the shim covers only the setter.

\Cref{fig:migration} summarizes the phased flow and the gate that ties it together.

\tikzfig{activity-migration}{The four-phase migration: the invariance gate is built
first (Phase~0), the new types and the analytic shim land with both APIs coexisting
(Phase~1), the gate is run against the migrated reader and the cone CM is pinned by hand
(Phase~2), and the fixtures are cut over---surfacing and fixing the pre-existing
poke-through---before the deprecated shim is retired (Phase~3).}{fig:migration}
